\documentclass[border=5mm]{standalone}
% ==============================================================================
% SETUP.TEX - CONFIGURACIÓN GLOBAL DEL PROYECTO
% ==============================================================================
% Este archivo centraliza todos los paquetes y estilos.
% Debe incluirse en el preámbulo del libro principal y de los archivos standalone.
% \PassOptionsToPackage{gray}{xcolor}
% --- 1. CODIFICACIÓN E IDIOMA ---
% fix-cm habilita las fuentes Computer Modern en tamaños arbitrarios (por debajo
% de 5 pt, entre otros). Debe cargarse antes que fontenc.
\usepackage{fix-cm}
\usepackage[utf8]{inputenc}
\usepackage[T1]{fontenc}
\usepackage[spanish, es-tabla]{babel}  
\usepackage{csquotes} % Recomendado para babel y citas

\usepackage{import}
% mode=image: Usa el PDF pre-compilado, nunca intenta recompilar.
% Debes compilar cada figura manualmente antes de compilar el libro.
\usepackage[mode=image, subpreambles=false]{standalone}

% --- 2. MATEMÁTICAS Y CIENCIAS ---
\usepackage{amsmath, amssymb, amsfonts, mathtools}
\usepackage{enumitem}

% LaTeX trae \DeclareMathSizes{5}{5}{5}{5}, así que a 5 pt (\tiny) los subíndices
% salen del mismo tamaño que la base: en las etiquetas de figuras el M de
% $\omega_M$ queda desproporcionado. Se restablece la proporción habitual.
\DeclareMathSizes{5}{5}{3.7}{3}

% --- 3. DISEÑO Y GEOMETRÍA --- 
% Unificamos los márgenes para todo el documento
\usepackage[left=2.5cm,right=2.5cm,top=3cm,bottom=3cm]{geometry}
\makeatletter
\ifstandalone % Evita que geometry dañe el recorte de las figuras bajándolas 3cm
  \@ifclasswith{standalone}{crop=false}{
    % Es un capítulo/sección: Dejamos que geometry actúe.
  }{
    % Es una figura: Neutralizamos geometry para que no las recorte mal.
    \geometry{pass}
  }
\fi
\makeatother
\usepackage{parskip} % Espaciado entre párrafos sin sangría
\usepackage{float}   % Para usar [H] en figura

% Ajustes de flotantes para evitar espacios verticales exagerados
\renewcommand{\topfraction}{0.95}      % Permite que las figuras ocupen hasta el 95% superior
\renewcommand{\bottomfraction}{0.95}   % Permite que las figuras ocupen hasta el 95% inferior
\renewcommand{\textfraction}{0.05}     % Permite hojas con tan solo un 5% de texto
\renewcommand{\floatpagefraction}{0.8} % Requiere que una página de solo figuras esté 80% llena
\setcounter{topnumber}{4}
\setcounter{bottomnumber}{4}
\setcounter{totalnumber}{10}

% --- 4. GRÁFICOS Y TABLAS ---
\usepackage{graphicx}
\usepackage{booktabs} % Tablas profesionales
\usepackage{caption}  % Control de captions y \captionof
\usepackage{tikz}
\usetikzlibrary{arrows.meta, calc, angles, quotes, babel, backgrounds}
\usepackage{pgfplots}
\usepgfplotslibrary{groupplots}
\usepgfplotslibrary{external}
\usepgfplotslibrary{patchplots}
\pgfplotsset{compat=1.18}
\usepackage[american, straightvoltages]{circuitikz}

% --- 5. COLORES Y CAJAS ---

\usepackage[table]{xcolor}
\usepackage{tcolorbox}

% Definición de colores corporativos/estilo
\definecolor{utnblue}{RGB}{0, 51, 102}
\definecolor{codegreen}{rgb}{0,0.6,0}
\definecolor{codegray}{rgb}{0.5,0.5,0.5}
\definecolor{codepurple}{rgb}{0.58,0,0.82}
\definecolor{backcolour}{rgb}{0.96,0.96,0.96}

% --- 6. ENTORNOS PERSONALIZADOS (CAJAS) ---

% Caja "Resultado" (Azul Corporativo) — Definiciones, fórmulas clave, resultados importantes.
% Uso: \begin{resultado}[Fórmula de Euler] ... \end{resultado}
\newtcolorbox{resultado}[1][]{
    colback=utnblue!5!white,
    colframe=utnblue,
    title=\textbf{#1},
    fonttitle=\bfseries,
    sharp corners,
    boxrule=0.5mm
}

% Caja "Nota" (Gris) — Advertencias, aplicaciones, contexto secundario.
% Uso: \begin{nota}[Cuidado con la calculadora] ... \end{nota}
% Uso: \begin{nota} ... \end{nota}  → título por defecto "Nota"
\newtcolorbox{nota}[1][Nota]{
    colback=codegray!5!white,
    colframe=codegray,
    title=\textbf{#1},
    fonttitle=\bfseries,
    sharp corners,
    boxrule=0.5mm
}

% Caja inline para resaltar un valor y enlazarlo con su otra aparición en una tabla
% o en una cuenta: mismo color, mismo valor.
% Uso: \marcavalor{$4$}  o  \marcavalor[codegreen]{$4$}
\newtcbox{\marcavalor}[1][utnblue]{
    on line,
    boxsep=1pt, left=2pt, right=2pt, top=1pt, bottom=1pt,
    colback=#1!10!white,
    colframe=#1,
    boxrule=0.4pt,
    arc=2pt
}

% --- 7. CÓDIGO FUENTE (LISTINGS) ---
\usepackage{listings}

\lstdefinestyle{miestilo}{
    backgroundcolor=\color{backcolour},   
    commentstyle=\color{codegreen},
    keywordstyle=\color{magenta},
    numberstyle=\tiny\color{codegray},
    stringstyle=\color{codepurple},
    basicstyle=\ttfamily\footnotesize,
    breakatwhitespace=false,         
    breaklines=true,                 
    captionpos=b,                    
    keepspaces=true,                 
    numbers=left,                    
    numbersep=5pt,                  
    showspaces=false,                
    showstringspaces=false,
    showtabs=false,                  
    tabsize=2,
    frame=single,                    
    rulecolor=\color{codegray}
}

\lstset{style=miestilo}
% Soporte para tildes y ñ en bloques de código
\lstset{literate={ñ}{{\~n}}1 {á}{{\'a}}1 {é}{{\'e}}1 {í}{{\'i}}1 {ó}{{\'o}}1 {ú}{{\'u}}1}

% Operadores matemáticos personalizados

% Vector en sentido discreto (lista finita de coordenadas): negrita + flecha.
% Uso: \vect{v}, \vect{e}_k, \vect{0}.
\newcommand{\vect}[1]{\vec{\mathbf{#1}}}

% Fecha de versión y comando para capítulos standalone
\providecommand{\verdate}{\the\year.\ifnum\month<10 0\fi\the\month.\ifnum\day<10 0\fi\the\day}
\newcommand{\chapterversion}{%
    \vspace{-1.5cm}\begin{flushright}\small\textit{\color{codegray}Versión~\verdate}\end{flushright}\vspace{0.5cm}%
}

% --- 8. HIPERVÍNCULOS (DEBE SER EL ÚLTIMO PAQUETE) ---
\usepackage[colorlinks=true, linkcolor=utnblue, urlcolor=utnblue, citecolor=utnblue]{hyperref}

% --- 9. REFERENCIAS CRUZADAS ENTRE CAPÍTULOS STANDALONE ---
% Cuando se compila un capítulo standalone, xr-hyper lee los .aux de los
% otros capítulos (si existen) para resolver \ref{} a labels externos.
% En la compilación del libro completo este bloque no se activa.
\ifstandalone
  \usepackage{xr-hyper}
  \makeatletter
  \newcommand{\xrchap}[1]{%
    \edef\xrc@self{\detokenize{Capitulo#1}}%
    \edef\xrc@job{\jobname}%
    \ifx\xrc@self\xrc@job\else
      \IfFileExists{../#1capitulo/Capitulo#1.aux}%
        {\externaldocument{../#1capitulo/Capitulo#1}}{}%
    \fi
  }
  \makeatother
  \xrchap{01}\xrchap{02}\xrchap{03}\xrchap{04}
  \xrchap{05}\xrchap{06}\xrchap{07}\xrchap{08}
  \xrchap{09}\xrchap{10}\xrchap{11}\xrchap{12}
  \xrchap{13}
\fi
% \PassOptionsToPackage{gray}{xcolor}

% Convolución de dos exponenciales causales (Ejemplo 1, sec. 4.3.2):
%   x(τ) = e^{-aτ} u(τ)        (a = 1)
%   h(τ) = e^{-bτ} u(τ)        (b = 2)
%
% Tres paneles:
% (a) Zona I, t₀ = -0.7: h(t₀-τ) está completamente a la izquierda de x(τ),
%     no hay solapamiento, y(t₀) = 0.
% (b) Zona II, t₀ = 1.0: solapamiento sobre [0, t₀], área = y(t₀).
% (c) Salida y(t) = (e^{-at} - e^{-bt})/(b-a) · u(t).
%     Máximo en t* = ln(b/a)/(b-a) = ln 2 ≈ 0.693, valor 0.25.

\begin{document}
\begin{tikzpicture}[>=Latex, font=\small]

\colorlet{xcol}{utnblue}
\colorlet{hcol}{red!80!black}

\pgfmathsetmacro{\tone}{-0.7}
\pgfmathsetmacro{\ttwo}{1.0}
\pgfmathsetmacro{\aval}{1.0}
\pgfmathsetmacro{\bval}{2.0}

% Mismo rango en los dos paneles del procedimiento (igual escala visual)
\pgfmathsetmacro{\xmn}{-2.6}
\pgfmathsetmacro{\xmx}{3.2}
\pgfmathsetmacro{\ymn}{-0.18}
\pgfmathsetmacro{\ymx}{1.25}

% =====================================================================
% PANEL (a): Zona I — t₀ < 0, sin solapamiento
% =====================================================================
\begin{axis}[
    name=panA,
    width=6cm, height=5cm,
    axis lines=center,
    xlabel={$\tau$},
    ylabel={},
    xlabel style={anchor=north west, at={(axis description cs:1,0)}},
    xmin=\xmn, xmax=\xmx,
    ymin=\ymn, ymax=\ymx,
    clip=true,
    axis line style={->, thick},
    xtick={\tone, 0},
    xticklabels={$t_0$, $0$},
    ytick={0,1},
    every tick label/.style={font=\tiny},
    title={\footnotesize (a)\; Zona~I:\; $t_0 < 0$},
]
    % x(τ) = e^{-aτ} u(τ)
    \addplot[xcol, very thick, domain=\xmn:0, samples=2] {0};
    \addplot[xcol, very thick, domain=0:\xmx, samples=120]
        {exp(-\aval*x)};
    \draw[xcol, very thick] (axis cs:0,0) -- (axis cs:0,1);
    \fill[xcol]  (axis cs:0,1) circle (2pt);
    \node[xcol, font=\tiny, above right] at (axis cs:0.05,1) {$x(\tau)$};

    % h(t₀-τ) con t₀ = -0.7: e^{b(τ-t₀)} para τ ≤ t₀
    \addplot[hcol, very thick, domain=\xmn:\tone, samples=120]
        {exp(\bval*(x-\tone))};
    \draw[hcol, very thick] (axis cs:\tone,0) -- (axis cs:\tone,1);
    \fill[hcol]  (axis cs:\tone,1) circle (2pt);
    \fill[white] (axis cs:\tone,0) circle (2pt);
    \draw[hcol]  (axis cs:\tone,0) circle (2pt);
    \addplot[hcol, very thick, domain=\tone:\xmx, samples=2] {0};
    \node[hcol, font=\tiny] at (axis cs:-1.7,0.55) {$h(t_0\!-\!\tau)$};

    % marcador t₀
    \draw[dashed, gray!60, thin] (axis cs:\tone,-0.16) -- (axis cs:\tone,1.1);
\end{axis}

% =====================================================================
% PANEL (b): Zona II — t₀ ≥ 0, solapamiento sobre [0, t₀]
% =====================================================================
\begin{axis}[
    name=panB,
    at={(panA.east)}, anchor=west, xshift=0.7cm,
    width=6cm, height=5cm,
    axis lines=center,
    xlabel={$\tau$},
    ylabel={},
    xlabel style={anchor=north west, at={(axis description cs:1,0)}},
    xmin=\xmn, xmax=\xmx,
    ymin=\ymn, ymax=\ymx,
    clip=true,
    axis line style={->, thick},
    xtick={0, \ttwo},
    xticklabels={$0$, $t_0$},
    ytick={0,1},
    every tick label/.style={font=\tiny},
    title={\footnotesize (b)\; Zona~II:\; $t_0 \geq 0$},
]
    % Solapamiento [0, t₀] sombreado
    \fill[gray!22] (axis cs:0,0) rectangle (axis cs:\ttwo,1);

    % x(τ)
    \addplot[xcol, very thick, domain=\xmn:0, samples=2] {0};
    \addplot[xcol, very thick, domain=0:\xmx, samples=120]
        {exp(-\aval*x)};
    \draw[xcol, very thick] (axis cs:0,0) -- (axis cs:0,1);
    \fill[xcol]  (axis cs:0,1) circle (2pt);
    \node[xcol, font=\tiny, above right] at (axis cs:0.05,1) {$x(\tau)$};

    % h(t₀-τ) con t₀ = 1
    \addplot[hcol, very thick, domain=\xmn:\ttwo, samples=120]
        {exp(\bval*(x-\ttwo))};
    \draw[hcol, very thick] (axis cs:\ttwo,0) -- (axis cs:\ttwo,1);
    \fill[hcol]  (axis cs:\ttwo,1) circle (2pt);
    \fill[white] (axis cs:\ttwo,0) circle (2pt);
    \draw[hcol]  (axis cs:\ttwo,0) circle (2pt);
    \addplot[hcol, very thick, domain=\ttwo:\xmx, samples=2] {0};
    \node[hcol, font=\tiny] at (axis cs:-0.6,0.55) {$h(t_0\!-\!\tau)$};

    % marcador t₀
    \draw[dashed, gray!60, thin] (axis cs:\ttwo,-0.16) -- (axis cs:\ttwo,1.1);

\end{axis}

% =====================================================================
% PANEL (c): Salida y(t)
% =====================================================================
\begin{axis}[
    name=panC,
    at={(panB.east)}, anchor=west, xshift=0.9cm,
    width=6cm, height=5cm,
    axis lines=center,
    xlabel={$t$},
    ylabel={$y(t)$},
    xlabel style={anchor=north west, at={(axis description cs:1,0)}},
    ylabel style={font=\small, rotate=0, anchor=south east, at={(axis description cs:0,1)}},
    xmin=-1.4, xmax=5.2,
    ymin=-0.05, ymax=0.32,
    clip=false,
    axis line style={->, thick},
    xtick={0, 0.6931},
    xticklabels={$0$, $t^{*}$},
    ytick={0, 0.25},
    yticklabels={$0$, $y_{\max}$},
    every tick label/.style={font=\tiny},
    title={\footnotesize (c)\; Salida $y(t)$},
]
    % Zona I: y = 0 para t < 0
    \addplot[domain=-1.3:0, samples=2, very thick, utnblue] {0};
    % Zona II: y = e^{-t} - e^{-2t} para t ≥ 0
    \addplot[domain=0:5.1, samples=160, very thick, utnblue]
        {exp(-x) - exp(-2*x)};
    \fill[utnblue] (axis cs:0,0) circle (1.8pt);
    % Máximo
    \fill[utnblue] (axis cs:0.6931, 0.25) circle (1.8pt);
    \draw[dashed, gray!60] (axis cs:0.6931, 0) -- (axis cs:0.6931, 0.25);
    \draw[dashed, gray!60] (axis cs:0, 0.25) -- (axis cs:0.6931, 0.25);
    % Etiquetas de zonas
    \node[font=\footnotesize, gray] at (axis cs:-0.7, 0.025) {Zona~I};
    \node[font=\footnotesize, gray] at (axis cs:1.5,  0.05)  {Zona~II};
    % Expresión analítica
    \node[font=\footnotesize, utnblue, align=left]
        at (axis cs:2.7, 0.20)
        {$\dfrac{e^{-at}-e^{-bt}}{b-a}$};
\end{axis}

\end{tikzpicture}
\end{document}
